\scanpage{16}
qui est donc un isom si $X/S$ est eff\textsuperscript{t} lisse et $S$ de car. nulle.

Je doute que ce soit tjrs un isom, \inserted{$\overline{\mathrm m}$} si $X$ est \uline{lisse et projectif} \inserted{(p.ex. une variété elliptique)} sur $S$ spectre d'un corps (alg. clos si on veut) de car $p>0$ ; \inserted{et $\underline V=\underline O_X$ avec stratification canonique.} Dans ce cas particulier, il n'est pas $\overline{\mathrm m}$ clair à priori que les $H^i_{\mathrm{str}}$ soient de dim finie, ni qu'ils soient nuls pour $i$ grand.
\uline{Situation à éclaircir !}
\margin{\figure{\begin{tikzpicture}[x=1cm,y=1cm]
\draw (1.6,6) ellipse (.53 and .37);
\draw (1.2,5.75) .. controls (-.2,4.6) and (-.7,4.8) .. (-.6,3.5) .. controls (-.6,2.9) and (-.2,2.6) .. (-.25,2.25);
\node[draw,ellipse,align=center,rotate=30,inner sep=3pt] at (0,1.65) {C'est en\\effet faux !};
\node[draw,align=center,rotate=30,inner sep=4pt] at (0,.15) {(\unsure{Calculs très}\\\unsure{amusants} pour\\$X/S$ \unsure{elliptique} ou\\abélien !)};
\draw (-.65,2.1)--(-.65,2.7);\draw (-.58,2.1)--(-.58,2.7);\draw (-.51,2.1)--(-.51,2.7);
\end{tikzpicture}}}

\uline{Remarque} L'intérêt de relier la coh. de De Rham à la cohomologie \deleted{\unsure{au sens}} $H_{\mathrm{str}}$ est que celle-ci, à son tour, s'interprète comme la cohomologie d'un ``site stratifiant'' convenable, dont une variante \unsure{ce|ns} \unsure{conduira} à la ``coh. $p$-adique en car. $p$''.
\margin{15}
